% Part: many-valued-logic
% Chapter: syntax-and-semantics
% Section: matrices

\documentclass[../../../include/open-logic-section]{subfiles}

\begin{document}

\olfileid{mvl}{syn}{mat}

\olsection{المصفوفات}

يتعين المنطق المتعدد القيم بأربعة أمور: لغته، ومجموعة قيم صدقه~${\OLId{latin-looped}{V}}$،
والمجموعة الجزئية المميزة من تلك القيم، ودوال الصدق لروابطه. ومجموع
هذه الأمور يسمى \emph{مصفوفة}.

\begin{defn}[المصفوفة]
\ollabel{defn:matrix}
\emph{مصفوفة} المنطق~$\Log L$ مؤلفة من:
\begin{enumerate}
\item مجموعة روابط هي لغة~$\Lang L$؛
\item مجموعة غير خالية~${\OLId{latin-looped}{V}} \neq \emptyset$ من قيم الصدق؛
\item مجموعة~${\OLId{latin-looped}{V}}^+ \subseteq {\OLId{latin-looped}{V}}$ من قيم الصدق المميزة؛
\item لكل رابط~$\star$ ذي~${\OLId{latin-ordinary}{n}}$ من المواضع في~$\Lang L$ دالة صدق
$\tf{\star} : {\OLId{latin-looped}{V}}^{\OLId{latin-ordinary}{n}} \to {\OLId{latin-looped}{V}}$. فإن كان~${\OLId{latin-ordinary}{n}} = {\OLMathNumeral{0}}$، لم تكن
$\tf{\star}$ إلا عنصرًا من~${\OLId{latin-looped}{V}}$.
\end{enumerate}
\end{defn}

\begin{ex}
أما مصفوفة المنطق الكلاسيكي~\LogCL{} فعناصرها:
\begin{enumerate}
  \item لغة القضايا القياسية~$\Lang L_{\OLMathNumeral{0}}$، وروابطها
  $\lfalse$ و$\lnot$ و$\land$ و$\lor$ و$\lif$.
  \item مجموعة قيم الصدق~${\OLId{latin-looped}{V}} = \{\True, \False\}$.
  \item القيمة المميزة وحدها هي~$\True$؛ أي~${\OLId{latin-looped}{V}}^+ = \{\True\}$.
  \item في الثابت~$\lfalse$ يكون~$\tf{\lfalse} = \False$؛ وأما
  دوال الصدق الباقية فمبينة في الجداول المعتادة (انظر
  \olref{fig:tf-CL}).
\end{enumerate}
\begin{figure}
  \begin{center}
      \begin{tabular}{c|c}
        $\tf{\lnot}$ & \\
        \hline
        $\True$ & $\False$ \\
        $\False$ & $\True$
      \end{tabular}
      \quad
      \begin{tabular}{c|cc}
        $\tf{\land}$ & $\True$ & $\False$ \\
        \hline
        $\True$ & $\True$ & $\False$ \\
        $\False$ & $\False$ & $\False$
      \end{tabular}
      \quad
      \begin{tabular}{c|cc}
        $\tf{\lor}$ & $\True$ & $\False$ \\
        \hline
        $\True$ & $\True$ & $\True$ \\
        $\False$ & $\True$ & $\False$
      \end{tabular}
      \quad
      \begin{tabular}{c|cc}
        $\tf{\lif}$ & $\True$ & $\False$ \\
        \hline
        $\True$ & $\True$ & $\False$ \\
        $\False$ & $\True$ & $\True$
      \end{tabular}
    \end{center}
    \caption{دوال الصدق في المنطق الكلاسيكي~$\LogCL$.}
    \ollabel{fig:tf-CL}
  \end{figure}
  \end{ex}

\end{document}
